\section{封装结构、海水空间与任意位姿场景}\label{sec:scene}
\subsection{真实材料域}
每个刚性装配件含 BFZI 线圈、导体截面、绝缘、灌封材料、
磁芯、屏蔽和独立参数化的封装盒外壳（可按工况选配）。
盒体的尺寸与壁厚不由 \(A_0\) 唯一决定，
盒体和线圈共享或分别指定合法的位姿。
实体导热区域记为 \(\Omega_i(t)\)，不允许不同材料区域非法重叠；
对气密封装，其海水域可定义为
\[
\Omega_s(t)=\mathbb R^3\setminus
\bigcup_p\overline{\Omega_{\mathrm{excluded},p}(t)}.
\]
这不是无限空间离散网格，而是几何定义；数值参考可以采取
人工截断域、无穷远边界或边界积分并进行收敛研究。

\subsection{几何材料查询与接口}
给每个组件局部坐标 \(\bm\xi_p=R_p^{\mathsf T}(\bm x-\bm c_p)\)，
材料边界可由解析几何、参数曲面或符号距离函数
\(d_i(\bm\xi_p;\bm g_i)\) 查询。
\textbf{SDF 是分区和求积辅助，不取代 BFZI 原始螺旋中心线。}
在非导电封装/海水界面，海水法向导电电流为零；
导电盒壁与海水相邻时必须使用真实电磁界面连续性，
不可套用绝缘条件。热学还必须区分固体内部导热、
接触传热及盒壁--海水热交换。

\subsection{材料范围和海水运动}
电磁材料可依赖 \(\bm x,T,\omega\)，包括
\(\sigma,\mu,\epsilon\)；热材料包括 \(\rho,c,k\)，
可分段均匀或在参考求解器允许时变参。
首版建议均匀各向同性海水、线性材料、
非导电非磁性封装；导电/磁性盒体、频散、非线性、
各向异性和任意空间变化是分阶段扩展，不能未经验证声称全覆盖。
世界坐标的海流 \(\bm v_s(\bm x,t)\)、重力和外界背景
不随组件的局部坐标旋转。任意\textbf{静态姿态}与高速运动不同：
若运动电动势、流固边界变化显著，应增加移动介质的电磁/热方程。
